% Shorter alternative to A01--A09 in abstract_introduction_revised.tex.
% Do not include this in addition to the full abstract. Retain Keywords and
% Introduction from the revised excerpt. Full numerical detail remains in
% the full-detail revision and should be retained in the Results section.
\textbf{\emph{Abstract} ---}
Traffic-signal decisions based on local conditions may overlook interactions between neighboring intersections and their consequences for network-wide delay. This study developed and evaluated two linked contributions: Pressure-Based Energy Optimization, which represents signal selection using an Ising-style objective combining traffic pressure, switching penalties, and neighboring-intersection coupling, and QAOA Global Optimization, which samples joint signal configurations using simulated parameterized quantum circuits.

Nine controllers were evaluated in microscopic SUMO simulations on $2\times2$ and $3\times3$ grids under 600-second saturated trials sweeping straight-movement probability and a 24-hour Seattle-derived demand profile. The energy-based controllers were particularly competitive under sustained demand and greater network interaction. In the original $3\times3$ saturated experiment at $\alpha=0.7$, QAOA completed a mean of 980 vehicles compared with 925 under Max-Pressure, with reported waiting times of 43.75~s and 47.16~s, respectively. Under full-day $3\times3$ Seattle-derived demand, QAOA's travel time was 92.96~s compared with 94.45~s for Classical Global. Comparisons across regimes showed trade-offs between completed-trip averages and throughput. Exact-energy diagnostics, repeated-run statistics, energy--traffic correlations, hyperparameter experiments, and synthetic angle, depolarizing, and readout perturbations showed that improved instantaneous energy quality did not consistently yield improved traffic performance.

A separate supplementary evaluation at $\alpha=0.7$ compared QAOA, global simulated annealing, and a frozen classical stochastic global controller using 20 optimizer seeds per controller and grid under fixed demand. Controller-difference analyses did not establish a QAOA traffic-performance advantage over the stochastic controller. Completed and unfinished populations, accumulated network-wide delay, and pre-entry delay were examined alongside completed-trip averages. Including pre-entry delay qualified the $2\times2$ improvement relative to global simulated annealing, while the combined delay measure favored QAOA on $3\times3$. A common-Hamiltonian replay compared recorded QAOA energies with three classical replays on each of 24,000 saved states. The $3\times3$ comparison passed the earlier descriptive screening tolerances, while $2\times2$ did not; neither result establishes distribution equivalence or matched computational cost.

The findings distinguish instantaneous surrogate-energy optimization from accumulated closed-loop traffic performance. Classical stochastic control reproduced some favorable outcomes observed with QAOA, supporting an interpretation that is not specific to quantum sampling. The experiments establish neither a quantum-specific nor a computational advantage, and the supplementary findings remain conditional on the tested fixed-demand scenarios.
